\documentclass[11pt,a4paper]{article}

\usepackage[utf8]{inputenc}
\usepackage[T1]{fontenc}
\usepackage[spanish,es-nodecimaldot]{babel}
\usepackage{geometry}
\usepackage{array}
\usepackage{longtable}
\usepackage{enumitem}
\usepackage{xcolor}
\usepackage{hyperref}

\geometry{margin=2.4cm}
\hypersetup{
  colorlinks=true,
  linkcolor=blue!55!black,
  urlcolor=blue!55!black
}

\setlength{\parskip}{0.55em}
\setlength{\parindent}{0pt}
\setlist[itemize]{topsep=0.25em,itemsep=0.18em,leftmargin=1.4em}
\setlist[enumerate]{topsep=0.25em,itemsep=0.18em,leftmargin=1.6em}

\newcommand{\sesion}[1]{\textbf{Sesión #1}}
\newcommand{\producto}{\textbf{Producto verificable:} }
\newcommand{\aidriven}{\textit{AI-driven}}

\title{\textbf{Plan de Taller de 4 Semanas}\\
\large Construcción paso a paso de una plataforma de sensado y monitoreo espectral}
\author{Propuesta de formación técnica}
\date{2026-05-21}

\begin{document}

\maketitle

\section*{Resumen ejecutivo}

Este documento propone un taller de 4 semanas para que un equipo técnico comprenda y pueda construir, paso a paso, una plataforma de sensado y monitoreo espectral. El taller no está planteado como una demostración de una plataforma ya terminada, sino como una experiencia guiada para entender su arquitectura, sus módulos, sus contratos de datos y su estilo de implementación.

El enfoque será \aidriven: se usará inteligencia artificial como apoyo para levantar requisitos, crear código base, documentar, revisar consistencia, preparar pruebas, depurar errores y explicar decisiones técnicas. No se contempla todavía IA embebida ni modelos corriendo en el sensor; la IA se usa como herramienta de trabajo durante el diseño y la construcción de la plataforma.

El stack se menciona solo como referencia práctica para construir el taller, no como el centro de la formación:

\begin{itemize}
  \item \textbf{Frontend:} Vite + React.
  \item \textbf{Backend:} Node.js con Express, pensado para operación multiusuario.
  \item \textbf{Tiempo real:} WebSocket o una capa equivalente para actualización en vivo.
  \item \textbf{Persistencia:} base de datos relacional, preferiblemente PostgreSQL.
  \item \textbf{Hardware:} cualquier SDR y cualquier embebido disponible para adquisición o simulación de datos.
  \item \textbf{Postprocesamiento:} Python, tratado como módulo externo y con baja profundidad, porque requiere más conocimientos de DSP y no hace parte del núcleo de la plataforma.
\end{itemize}

\section*{Propósito del taller}

El propósito es que los participantes puedan entender cómo se diseña y se construye una plataforma que conecta una interfaz web, un backend, una base de datos, sensores o simuladores, hardware SDR/embebido y un módulo externo de análisis.

Al terminar, el equipo debe poder:

\begin{itemize}
  \item Explicar la arquitectura general de la plataforma.
  \item Identificar responsabilidades de frontend, backend, base de datos, sensor, embebido y postprocesamiento.
  \item Entender bases mínimas de RF y SDR para que el sensor no sea una caja negra.
  \item Levantar un entorno base usando herramientas web modernas.
  \item Entender el flujo de datos desde el SDR o simulador hasta la interfaz web.
  \item Diseñar endpoints, eventos en tiempo real y modelos de datos básicos.
  \item Aplicar un estilo de código ordenado y mantenible.
  \item Usar IA de forma responsable para acelerar requisitos, código, pruebas y documentación.
  \item Presentar una demo funcional o semi-funcional de punta a punta.
\end{itemize}

\section*{Audiencia objetivo}

El taller está dirigido a un equipo técnico pequeño:

\begin{itemize}
  \item Un ingeniero de sistemas.
  \item Dos o más tecnólogos o perfiles técnicos equivalentes.
\end{itemize}

No se asume dominio previo de radiofrecuencia, DSP o procesamiento avanzado de señales. Sí se recomienda que los participantes tengan bases de programación, uso de terminal, Git, HTTP y desarrollo web básico.

\section*{Enfoque pedagógico}

El taller se centra en aprender a construir y razonar la plataforma, no solo en verla funcionando. Cada semana combina explicación arquitectónica, revisión de estilo de código, laboratorios guiados y entregables verificables.

La profundidad del código será selectiva:

\begin{itemize}
  \item Se revisarán estructuras, patrones y responsabilidades.
  \item Se construirán esqueletos funcionales.
  \item Se explicarán contratos de API, eventos, modelos y flujos.
  \item No se hará una explicación exhaustiva de cada línea.
  \item Se darán bases pequeñas de RF/SDR para entender antenas, frecuencia, potencia, espectro y representación de datos.
  \item Se evitará profundizar en DSP avanzado.
  \item Se usará IA para generar borradores, revisar consistencia y documentar decisiones.
\end{itemize}

\section*{Rol de la inteligencia artificial}

El enfoque \aidriven{} busca que los participantes aprendan a usar IA como copiloto técnico durante la creación de requisitos, código, documentación, pruebas y material de soporte. La IA no reemplaza el criterio del equipo; ayuda a acelerar ciclos de diseño, documentación y validación.

Este taller no propone todavía IA dentro del embebido, del SDR o del sensor. El hardware se mantiene como fuente de adquisición de datos o como punto de envío hacia la plataforma.

\subsection*{Usos esperados de IA}

\begin{itemize}
  \item Convertir requisitos en historias técnicas.
  \item Proponer diagramas de arquitectura.
  \item Generar esqueletos de frontend y backend.
  \item Proponer modelos de datos iniciales.
  \item Revisar consistencia entre endpoints, DTOs y componentes.
  \item Sugerir pruebas manuales y automáticas.
  \item Explicar errores de compilación o ejecución.
  \item Documentar decisiones técnicas.
  \item Preparar guías de operación y handover.
  \item Generar ejemplos de prompts para mejorar código, requisitos y documentación.
\end{itemize}

\subsection*{Reglas de uso de IA durante el taller}

\begin{itemize}
  \item Todo código generado por IA debe ser revisado por una persona.
  \item La IA debe recibir contexto claro: objetivo, archivo, restricción y salida esperada.
  \item No se aceptan respuestas generadas sin validar contra el flujo real de la plataforma.
  \item Se debe pedir a la IA que explique supuestos y riesgos.
  \item Se debe conservar un registro breve de decisiones importantes.
  \item No se diseña ningún componente de IA embebida durante este taller.
\end{itemize}

\section*{Requisitos de hardware}

El taller debe ser flexible frente al hardware disponible. No se amarra a una referencia específica de SDR o embebido.

\begin{itemize}
  \item Cualquier SDR disponible, por ejemplo RTL-SDR, HackRF, USRP, LimeSDR, SDRplay u otro compatible.
  \item Cualquier embebido o equipo de borde disponible, por ejemplo Raspberry Pi, Jetson, BeagleBone, mini-PC industrial u otro equipo capaz de capturar, preparar o enviar datos.
  \item Antena compatible con el rango de prueba.
  \item Computadores para los participantes.
  \item Red local o acceso a internet para conectar frontend, backend y equipo de adquisición.
  \item En caso de no tener hardware listo, se usará simulador de datos.
\end{itemize}

\section*{Requisitos de software}

\begin{itemize}
  \item Node.js LTS.
  \item npm, pnpm o yarn.
  \item Vite + React para el frontend.
  \item Node.js + Express para el backend.
  \item TypeScript recomendado para frontend y backend.
  \item PostgreSQL local o en contenedor.
  \item Docker opcional para simplificar servicios.
  \item Git.
  \item Cliente HTTP: Postman, Insomnia, Thunder Client o equivalente.
  \item Python 3.10+ para postprocesamiento o simuladores auxiliares.
  \item Drivers o librerías necesarias para el SDR elegido.
  \item Acceso a una herramienta de IA para asistencia técnica.
\end{itemize}

\section*{Arquitectura de referencia}

La plataforma se entiende como un sistema por capas:

\begin{enumerate}
  \item \textbf{Capa de adquisición:} SDR y embebido capturan o simulan datos de campo.
  \item \textbf{Módulo de adquisición:} script o proceso simple que prepara y envía estado, ubicación, espectro y audio cuando aplique.
  \item \textbf{Backend Express:} coordina usuarios, sensores, configuraciones, campañas, datos, alertas y reportes.
  \item \textbf{Base de datos:} guarda configuración, histórico, campañas, mediciones y resultados.
  \item \textbf{Tiempo real:} WebSocket transmite estados y datos recientes al frontend.
  \item \textbf{Frontend React:} muestra mapas, sensores, monitoreo, campañas, alertas y resultados.
  \item \textbf{Postprocesamiento Python:} módulo externo para análisis especializado, tratado como integración puntual.
\end{enumerate}

\section*{Estilo de código esperado}

El taller debe enseñar un estilo de código ordenado, legible y extensible. La meta es que el equipo pueda continuar la plataforma después del taller.

\subsection*{Backend}

\begin{itemize}
  \item Separar rutas, controladores, servicios, repositorios y modelos.
  \item Mantener controladores delgados: reciben, validan y delegan.
  \item Concentrar reglas de negocio en servicios.
  \item Usar DTOs o schemas para entradas y salidas.
  \item Centralizar configuración de entorno.
  \item Manejar errores de forma consistente.
  \item Diseñar endpoints con contratos claros antes de implementar.
  \item Pensar desde el inicio en multiusuario: sesiones, roles, permisos y auditoría.
\end{itemize}

\subsection*{Frontend}

\begin{itemize}
  \item Separar páginas, componentes, hooks, servicios API y estado global.
  \item Mantener componentes de UI simples y reutilizables.
  \item Evitar lógica de negocio pesada dentro de componentes visuales.
  \item Definir estados de carga, error, vacío y éxito.
  \item Documentar flujos de usuario antes de construir pantallas.
  \item Priorizar una interfaz clara para operación técnica.
\end{itemize}

\subsection*{Integraciones}

\begin{itemize}
  \item Definir contratos de datos antes de conectar hardware.
  \item Simular datos antes de depender del SDR real.
  \item Versionar cambios en endpoints y eventos.
  \item Registrar logs útiles para diagnóstico.
  \item Documentar supuestos de frecuencia, resolución, timestamp y unidades.
\end{itemize}

\section*{Duración sugerida}

Se propone una duración de 4 semanas con 3 sesiones por semana, cada una de 2 horas. Esto da 24 horas sincrónicas. Se recomienda complementar con 2 a 4 horas semanales de práctica autónoma.

\begin{itemize}
  \item \textbf{Semanas:} 4.
  \item \textbf{Sesiones por semana:} 3.
  \item \textbf{Duración por sesión:} 2 horas.
  \item \textbf{Total sincrónico:} 24 horas.
  \item \textbf{Modalidad:} presencial, remota o mixta.
  \item \textbf{Formato:} explicación breve, laboratorio guiado, discusión de arquitectura y cierre con entregable.
\end{itemize}

\section*{Cronograma detallado}

\begin{longtable}{|p{0.13\textwidth}|p{0.24\textwidth}|p{0.35\textwidth}|p{0.18\textwidth}|}
\hline
\textbf{Semana} & \textbf{Sesión} & \textbf{Contenido y actividad} & \textbf{Entregable} \\
\hline
\endhead

1 & \sesion{1}: visión, RF y SDR básico &
Presentación del problema y bases mínimas de radio: espectro electromagnético, RF, antenas, potencia, frecuencia, ancho de banda y qué entrega un SDR. El objetivo es que el sensor no se trate como una caja negra. &
Glosario RF/SDR básico y primer mapa del flujo de datos. \\
\hline

1 & \sesion{2}: arquitectura y estilo de proyecto &
Actores del sistema, flujo desde sensor hasta interfaz, módulos principales y responsabilidades. Mención breve del stack de referencia. Uso de IA para convertir el problema en arquitectura, requisitos iniciales y lista de capacidades. &
Mapa inicial de arquitectura, estructura base propuesta y guía corta de estilo. \\
\hline

1 & \sesion{3}: entorno y primer flujo mínimo &
Preparación del ambiente. Creación del frontend base, backend base, endpoint de salud y cliente de prueba. Uso de IA para generar checklist de instalación, diagnóstico y documentación inicial. &
Frontend y backend ejecutándose localmente con primer endpoint probado. \\
\hline

2 & \sesion{4}: modelo de datos &
Diseño de entidades: usuario, sensor, antena, configuración, medición, campaña, alerta y reporte. Relación entre datos operativos e históricos. Discusión de PostgreSQL y migraciones. &
Modelo entidad-relación inicial y lista de tablas mínimas. \\
\hline

2 & \sesion{5}: API Express multiusuario &
Diseño de API REST. Separación de rutas, controladores, servicios y repositorios. Sesiones, roles y permisos a nivel conceptual. Construcción de endpoints base para sensores y campañas. &
Contrato API inicial y backend modular. \\
\hline

2 & \sesion{6}: recepción de datos y tiempo real &
Diseño del contrato del sensor: estado, GPS, espectro y audio opcional. Ingesta por HTTP y publicación por WebSocket. Simulación de datos antes de usar hardware real. &
Flujo backend capaz de recibir datos simulados y emitir eventos en vivo. \\
\hline

3 & \sesion{7}: arquitectura frontend &
Estructura del frontend con la herramienta de referencia: rutas, layout, menú, páginas, componentes, servicios API y manejo de estado. Definición de pantallas principales: inicio, dispositivos, monitoreo, campañas, alertas y configuración. &
Mapa de navegación y esqueleto de pantallas. \\
\hline

3 & \sesion{8}: visualización e interacción &
Conexión con API. Estados de carga, error y datos vacíos. Visualización de sensores, estado general, mediciones y campañas. Uso de IA para revisar consistencia entre flujo de usuario y componentes. &
Frontend conectado al backend con datos de prueba. \\
\hline

3 & \sesion{9}: monitoreo en vivo &
Integración de WebSocket. Actualización de datos en vivo. Vista de monitoreo, selección de sensor, parámetros de medición y lectura de datos recientes. Discusión de estilo visual operativo. &
Pantalla de monitoreo funcional con datos simulados o reales. \\
\hline

4 & \sesion{10}: SDR, embebido y adquisición &
Conexión conceptual con cualquier SDR y cualquier embebido. Diseño del flujo de adquisición: captura o simulación, empaquetado de datos, envío al backend y manejo de errores. Si no hay hardware, se usa simulador. &
Contrato de adquisición y prueba de envío desde hardware o simulador. \\
\hline

4 & \sesion{11}: reportes y postprocesamiento Python &
Explicación liviana del rol de Python: análisis externo, generación de resultados y contrato con backend. Se recalca que DSP no es el centro del taller. Diseño de flujo para reporte o resultado de campaña. &
Contrato de integración con Python y flujo de reporte definido. \\
\hline

4 & \sesion{12}: integración final y handover &
Prueba punta a punta: sensor o simulador, backend, base de datos, WebSocket, frontend y reporte básico. Revisión de riesgos, documentación mínima, tareas pendientes y demo final. &
Demo final, guía de ejecución y backlog de continuidad. \\
\hline
\end{longtable}

\section*{Plan por semanas}

\subsection*{Semana 1: fundamentos RF/SDR, arquitectura y entorno}

\textbf{Objetivo:} que el equipo entienda qué se va a construir, cómo fluye una medición desde el mundo RF hasta la interfaz y por qué la plataforma se separa en módulos.

\textbf{Temas principales:}

\begin{itemize}
  \item Propósito de la plataforma.
  \item Bases del espectro electromagnético y radiofrecuencia.
  \item Antenas: rol, bandas, ganancia, polarización y relación con la medición.
  \item SDR: qué mide, qué no mide y cómo entrega datos.
  \item Representación de datos: frecuencia, potencia, muestras, FFT, resolución y ancho de banda a nivel introductorio.
  \item Diferencia entre señal física, dato capturado y dato visualizado.
  \item Flujo desde SDR/embebido hasta frontend.
  \item Diferencia entre demostración, prototipo y plataforma multiusuario.
  \item Arquitectura por capas.
  \item Stack tecnológico como referencia, sin convertirlo en el foco del curso.
  \item Reglas de estilo de código.
  \item Uso de IA para requisitos, arquitectura, documentación y generación de esqueletos.
\end{itemize}

\textbf{Laboratorios:}

\begin{itemize}
  \item Construir un glosario mínimo de RF/SDR.
  \item Dibujar el recorrido de una señal desde antena hasta pantalla.
  \item Crear diagrama de arquitectura asistido por IA.
  \item Crear estructura base del proyecto.
  \item Levantar el frontend con la herramienta de referencia.
  \item Levantar el backend con la herramienta de referencia.
  \item Crear endpoint de salud.
  \item Probar backend desde cliente HTTP.
\end{itemize}

\producto repositorio base ejecutándose localmente, glosario RF/SDR, diagrama inicial y guía de estilo.

\subsection*{Semana 2: backend, datos y tiempo real}

\textbf{Objetivo:} construir la columna vertebral de la plataforma: API, modelos, persistencia, recepción de datos y eventos en vivo.

\textbf{Temas principales:}

\begin{itemize}
  \item Diseño de datos.
  \item Backend modular con Express.
  \item Multiusuario: usuarios, roles, permisos y sesiones.
  \item Gestión de sensores y antenas.
  \item Gestión de campañas.
  \item Recepción de datos de sensor.
  \item Eventos en tiempo real.
  \item Pruebas básicas de endpoints.
  \item Uso de IA para generar código base, revisar contratos y documentar endpoints.
\end{itemize}

\textbf{Laboratorios:}

\begin{itemize}
  \item Diseñar tablas principales.
  \item Crear endpoints base.
  \item Crear un simulador simple de sensor.
  \item Recibir datos de estado, ubicación y espectro.
  \item Publicar eventos por WebSocket.
  \item Documentar contratos de API.
\end{itemize}

\producto backend funcional con base de datos, endpoints principales y flujo de datos simulado.

\subsection*{Semana 3: frontend e integración}

\textbf{Objetivo:} construir la experiencia web que permita operar, consultar y visualizar la plataforma.

\textbf{Temas principales:}

\begin{itemize}
  \item Arquitectura de frontend con la herramienta definida para el taller.
  \item Rutas, layout y menú.
  \item Servicios API.
  \item Manejo de estado.
  \item Pantallas de inicio, dispositivos, monitoreo, campañas, alertas y configuración.
  \item Visualización de datos en vivo.
  \item Estados de interfaz: cargando, error, vacío, éxito.
  \item Uso de IA para crear componentes base, revisar estados de interfaz y documentar decisiones.
\end{itemize}

\textbf{Laboratorios:}

\begin{itemize}
  \item Crear layout principal.
  \item Crear servicios de consumo de API.
  \item Mostrar sensores y campañas.
  \item Conectar vista de monitoreo a WebSocket.
  \item Revisar consistencia visual y técnica con IA.
\end{itemize}

\producto frontend conectado al backend, con flujo de monitoreo visible y datos simulados o reales.

\subsection*{Semana 4: hardware, postprocesamiento liviano, pruebas y entrega}

\textbf{Objetivo:} cerrar un flujo completo, integrar hardware o simulador y preparar una demo final entendible.

\textbf{Temas principales:}

\begin{itemize}
  \item Conexión conceptual con cualquier SDR.
  \item Uso de cualquier embebido como nodo de adquisición.
  \item Módulo simple de captura o envío de datos.
  \item Contrato entre adquisición y backend.
  \item Postprocesamiento Python como módulo externo.
  \item Reportes y resultados.
  \item Pruebas end-to-end.
  \item Documentación de operación.
  \item Handover técnico.
\end{itemize}

\textbf{Laboratorios:}

\begin{itemize}
  \item Enviar datos desde SDR/embebido o simulador.
  \item Validar recepción en backend.
  \item Ver actualización en frontend.
  \item Ejecutar flujo de campaña o monitoreo.
  \item Generar resultado o reporte básico.
  \item Preparar demo final.
\end{itemize}

\producto demo de punta a punta, documentación mínima, lista de riesgos y backlog de continuidad.

\section*{Alcance del postprocesamiento Python}

El postprocesamiento se menciona y se integra de forma controlada, pero no será el centro del taller. La razón es que el análisis de señales y DSP requiere conocimientos especializados adicionales.

Durante el taller se verá:

\begin{itemize}
  \item Qué recibe Python desde el backend.
  \item Qué devuelve Python al backend.
  \item Cómo se invoca el módulo externo.
  \item Cómo se presenta un resultado en la interfaz.
  \item Qué límites tiene el análisis si no se entra en DSP.
\end{itemize}

No se profundizará en:

\begin{itemize}
  \item Teoría avanzada de FFT.
  \item Detección robusta de emisiones.
  \item Demodulación avanzada.
  \item Calibración de RF.
  \item Modelos matemáticos de cumplimiento.
\end{itemize}

\section*{Estructura sugerida del repositorio}

La estructura exacta puede cambiar, pero el taller debe dejar claro cómo separar responsabilidades.

\begin{verbatim}
platform/
  frontend/
    src/
      components/
      pages/
      hooks/
      services/
      state/
      styles/
  backend/
    src/
      routes/
      controllers/
      services/
      repositories/
      models/
      websocket/
      config/
  acquisition-client/
    src/
      capture/
      transport/
      config/
  postprocessing/
    src/
      adapters/
      reports/
  docs/
    architecture/
    api/
    runbooks/
\end{verbatim}

\section*{Entregables finales}

Al finalizar las 4 semanas se espera entregar:

\begin{itemize}
  \item Diagrama de arquitectura.
  \item Guía de instalación y ejecución.
  \item Repositorio base o estructura de referencia.
  \item Backend Express con endpoints principales.
  \item Frontend con pantallas principales usando la herramienta definida para el taller.
  \item Flujo de datos desde simulador, SDR o embebido.
  \item WebSocket o mecanismo equivalente de datos en vivo.
  \item Contrato básico para postprocesamiento Python.
  \item Demo final.
  \item Backlog de mejoras.
\end{itemize}

\section*{Criterios de éxito}

El taller será exitoso si los participantes pueden:

\begin{itemize}
  \item Explicar la plataforma sin depender del instructor.
  \item Identificar dónde vive cada responsabilidad técnica.
  \item Crear un nuevo endpoint siguiendo el estilo del backend.
  \item Crear una nueva pantalla siguiendo el estilo del frontend.
  \item Conectar una fuente de datos simulada o real.
  \item Diagnosticar errores básicos de API, base de datos y WebSocket.
  \item Usar IA para acelerar tareas sin copiar ciegamente.
  \item Presentar una demo técnica coherente.
\end{itemize}

\section*{Riesgos y mitigaciones}

\begin{itemize}
  \item \textbf{Hardware no disponible o inestable:} usar simulador como plan base y dejar SDR/embebido como integración progresiva de adquisición.
  \item \textbf{Exceso de profundidad en código:} mantener el foco en arquitectura, contratos y estilo.
  \item \textbf{Dificultad de DSP:} tratar Python como caja externa con contrato claro.
  \item \textbf{Dependencia excesiva de IA:} exigir revisión humana y explicación de decisiones.
  \item \textbf{Diferencias de nivel técnico:} trabajar por parejas y dejar tareas de refuerzo.
  \item \textbf{Alcance demasiado grande:} priorizar flujo mínimo de punta a punta antes de funcionalidades avanzadas.
\end{itemize}

\section*{Agenda sugerida de cada sesión}

Cada sesión de 2 horas puede seguir este formato:

\begin{itemize}
  \item 15 minutos: objetivo y repaso del flujo.
  \item 25 minutos: explicación arquitectónica.
  \item 45 minutos: laboratorio guiado.
  \item 20 minutos: revisión con IA y discusión de estilo.
  \item 15 minutos: cierre, dudas y tarea corta.
\end{itemize}

\section*{Material mínimo para el instructor}

\begin{itemize}
  \item Diagrama de arquitectura base.
  \item Repositorio inicial o plantilla.
  \item Prompts guía para IA.
  \item Colección de endpoints de prueba.
  \item Simulador de sensor.
  \item Datos de ejemplo.
  \item Checklist de instalación.
  \item Checklist de demo final.
\end{itemize}

\section*{Conclusión}

Este taller de 4 semanas busca que el equipo aprenda a construir la plataforma paso a paso, entendiendo primero bases mínimas de RF/SDR y la arquitectura, y luego bajando a implementaciones guiadas. El stack de referencia permite cubrir el flujo completo sin convertir el taller en una clase extensa de DSP ni en una revisión exhaustiva de código.

El resultado esperado es un equipo capaz de explicar, ejecutar, modificar y continuar la plataforma con criterio técnico, apoyo de IA para producir código/documentación/requisitos, y una base de trabajo ordenada.

\end{document}